\chapter*{Resumo}
\noindent
A depuração de sistemas embarcados exige instrumentos capazes de capturar
sinais digitais com precisão temporal, mas o custo de analisadores lógicos
comerciais restringe seu acesso em contextos educacionais e de pesquisa.
Um dispositivo de aquisição de dados (DAQ) de baixo custo baseado no
microcontrolador RP2350, compatível com o ecossistema Sigrok/PulseView,
delega a amostragem ao subsistema de E/S programável (PIO) e a
movimentação dos dados ao controlador de acesso direto à memória (DMA);
o ganho dessa arquitetura sobre a alternativa convencional, a leitura dos
pinos de entrada e saída de propósito geral (GPIO) pela própria CPU, no
entanto, não havia sido quantificado. Este trabalho implementa uma segunda
versão do firmware, com aquisição integralmente por software em núcleo
único, mantendo
idênticos o protocolo de comunicação com o \emph{host}, a pilha USB e a
codificação RLE, e compara experimentalmente as duas arquiteturas quanto à
taxa máxima de amostragem sustentada, ao \emph{jitter} de amostragem, à
perda de amostras em captura contínua e à ocupação de CPU. As medições
usaram uma onda quadrada de referência, um analisador lógico comercial
como referência externa e 30 repetições por condição, com comparação pelo
teste de Mann-Whitney e correção de Holm-Bonferroni. A arquitetura PIO+DMA
sustentou 3,66\,MS/s com contagem de bordas exata, \emph{jitter} igual ao
limite de quantização e nenhuma amostra perdida em captura contínua até
3,54\,MS/s, abortando a captura de forma sinalizada fora dessa faixa, com
ocupação de CPU entre 0,08\,\% e 83\,\%, proporcional à taxa. A leitura
por CPU ocupou 100\,\% do processador e, embora o custo de seu laço
permitisse cerca de 4,29\,MS/s, falhou em silêncio a partir de
$\approx 1{,}1$\,MS/s: o laço para enquanto a mesma CPU transmite os
dados, o que produziu bordas a mais nas capturas com disparo,
\emph{jitter} de 106 a 400\,ns acima da quantização e perda de 24\,\% a
43\,\% das amostras em captura contínua. Conclui-se que a leitura por CPU em
núcleo único, com o mesmo núcleo amostrando e transmitindo, basta para
capturas curtas em taxas de até centenas de kS/s, enquanto o
par PIO+DMA é indispensável em taxas maiores, em captura contínua ou
quando a CPU precisa ficar livre. A arquitetura determina, além da taxa, o
custo de detectar a falha: no par PIO+DMA o transbordo da fila é um evento de
hardware que a CPU observa fora do caminho crítico, ao passo que, na leitura
por software, detectar as pausas consome ciclos do próprio laço de
amostragem.

\vspace{5mm}
\noindent\textbf{Palavras-chave:} Aquisição de dados. Microcontroladores. RP2350. PIO. DMA.
